\clearpage
\begin{tcolorbox}
[title=Prompt for Extracting the Final Answer, colback=white, colframe=black!75!white, breakable]
\begin{spacing}{1.3}
\refstepcounter{tcbtable}
\label{prompt:correct}

% You are an ``Answer-Level Trajectory Recorder''.

% The user will provide a QA pair. Your goal is to analyze ONLY the Answer (A) and extract the full trajectory of ANSWER-LEVEL STATES for the TARGET quantity asked by the question.

% This task is about analyzing how the model’s ANSWER changes over time.

% --------------------------------\\
% PHASE 1: ANALYSIS (FREE-FORM)\\
% --------------------------------\\
% First, read the Question (Q) ONLY to identify the TARGET quantity
% (i.e., what the question ultimately asks the model to answer).\\

% Then analyze the Answer (A) and reason about how the model treats different values or expressions as the ANSWER to that TARGET over time.

% In this phase, you should:\\
% - Distinguish answer-level commitments from intermediate computations.\\
% - Identify moments where the model:\\
%   -proposes a candidate final answer\\
%   -temporarily commits to an answer\\
%   -revises or replaces a previously stated answer\\
%   -re-confirms or restates the same answer again\\

% --------------------------------\\
% PHASE 2: FINAL OUTPUT (STRICT)\\
% --------------------------------\\
% After the analysis, output the ANSWER-LEVEL TRAJECTORY in chronological order.

% ========================\\
% WHAT COUNTS AS AN ANSWER-LEVEL STATE\\
% ========================\\

% An answer-level state MUST satisfy BOTH conditions:

% (1) Target Alignment:\\
% - It refers to the TARGET quantity asked by Q
%   (e.g., “the value of f(-2)+f(-1)+f(0)”, “the final common fraction”, “the correct option”).

% (2) Answer Commitment:\\
% - The model presents it as an answer, conclusion, or result for the TARGET,\\
%   including:\\
%   -``the answer is …''\\
%   -``the value is …''\\
%   -``thus … equals …''\\
%   -``final answer: …''\\
%   -an explicit restatement or confirmation of a previously given answer

% IMPORTANT:\\
% A statement can count as an answer-level state EVEN IF it later turns out to be wrong.\\

% ========================\\
% WHAT MUST BE FILTERED OUT\\
% ========================\\

% Do NOT extract:\\
% - Sub-results used only to compute the TARGET
%   (e.g., f(-2)=2, f(-1)=5/3, f(0)=1)\\
% - Pure intermediate calculations or simplification steps\\
% - Values that are not presented as the answer to the TARGET\\

% Even if such items are boxed or emphasized, they are NOT answer-level states unless they are treated as the answer to the TARGET.

% ========================\\
% TRAJECTORY RULES\\
% ========================\\

% - Extract answer-level states strictly in the order they appear in A.\\
% - Do NOT merge, deduplicate, or normalize answers.\\
% - If the same answer is stated or confirmed multiple times, record EACH occurrence as a separate state.\\
% - If multiple different candidate final answers appear, record ALL of them.\\
% - Do NOT invent answer states that do not explicitly appear in A.\\

% ========================\\
% FINAL OUTPUT FORMAT (MANDATORY)\\
% ========================\\

% Your final output must contain ONLY boxed answer-level states, one per line.\\
% No explanations, no headers, no numbering, no blank lines.\\

% Each line must follow exactly this format:\\
% \\boxed{}
% \boxed{...}
You are an \textbf{``Answer-Level Trajectory Recorder''}.  

The user will provide a QA pair. Your goal is to analyze \textbf{ONLY the Answer (A)} and extract the full trajectory of \textbf{ANSWER-LEVEL STATES} for the \textbf{TARGET} quantity asked by the question.  
This task analyzes how the model’s ANSWER changes over time.

\textbf{Phase 1: Analysis (Free-Form)}  
- First, read the Question (Q) ONLY to identify the \textbf{TARGET} quantity (what the question ultimately asks the model to answer).  
- Then analyze the Answer (A) and reason about how the model treats different values or expressions as the ANSWER to that TARGET over time.  
- In this phase, distinguish answer-level commitments from intermediate computations, and identify moments where the model:  
\begin{itemize}[leftmargin=*, label={--}]
    \item proposes a candidate final answer  
    \item temporarily commits to an answer  
    \item revises or replaces a previously stated answer  
    \item re-confirms or restates the same answer
\end{itemize}

\textbf{Phase 2: Final Output (Strict)}  
- After the analysis, output the \textbf{ANSWER-LEVEL TRAJECTORY} in chronological order.

\textbf{What Counts as an Answer-Level State}  
An answer-level state MUST satisfy BOTH conditions:

\begin{enumerate}[leftmargin=*, label=(\arabic*)]
    \item \textbf{Target Alignment:} Refers to the TARGET quantity asked by Q (e.g., “the value of f(-2)+f(-1)+f(0)”, “the final common fraction”, “the correct option”).  
    \item \textbf{Answer Commitment:} The model presents it as an answer, conclusion, or result for the TARGET, including:  
    \begin{itemize}[leftmargin=*, label={--}]
        \item ``the answer is …''  
        \item ``final answer: …''  
        \item an explicit restatement or confirmation of a previously given answer
    \end{itemize}  
    \textit{Important:} A statement counts even if it later turns out to be wrong.
\end{enumerate}

\textbf{What Must Be Filtered Out}  
Do NOT extract:  
\begin{itemize}[leftmargin=*, label={--}]
    \item Sub-results used only to compute the TARGET (e.g., f(-2)=2, f(-1)=5/3, f(0)=1)  
    \item Pure intermediate calculations or simplification steps  
    \item Values not presented as the answer to the TARGET
\end{itemize}  
Even if boxed or emphasized, they are NOT answer-level states unless treated as the answer.

\textbf{Trajectory Rules}  
\begin{itemize}[leftmargin=*, label={--}]
    \item Extract answer-level states strictly in the order they appear in A.  
    \item Do NOT merge, deduplicate, or normalize answers.  
    \item If the same answer is stated or confirmed multiple times, record EACH occurrence separately.  
    \item If multiple different candidate final answers appear, record ALL of them.  
    % \item Do NOT invent answer states that do not explicitly appear in A.
\end{itemize}

\textbf{Final Output Format (Mandatory)}  
Your output must contain ONLY \textbf{boxed answer-level states}, one per line. No explanations, headers, numbering, or blank lines.  

\begin{verbatim}
\boxed{}
\end{verbatim}
\end{spacing}
\end{tcolorbox}